\section{Технический проект}

\subsection{Структура диска}

Диск имеет следующую структуру:

\begin{enumerate}
\item Первый сектор(512 байт) диска занимает главная загрузочная запись(MBR). Главная загрузочная запись содержит следующую информацию:

	\begin{itemize}
	\item код загрузчика(446 байт);
	\item таблица разделов(4 элемента по 16 байт = 64 байта);
	\item сигнатура(2 байта). Всегда имеет значение 21930($55AA_{16}$).
	\end{itemize}
	
Каждая запись раздела из таблицы разделов содержит следующую информацию:

	\begin{enumerate}
	\item Признак активности раздела(1 байт), который может принимать только 2 значения: 0(свободный раздел) или 128($80_{16}$)(раздел занят).
	\item Адрес начала раздела(3 байта) в виде номеров цилиндра, головки, сектора.
	\item Код типа файловой системы(1 байт). 0 -- пустая запись. 12($C_{16}$) -- FAT-32. 148($94_{16}$) -- ISO-9660(CDFS).
	\item Адрес конца раздела(3 байта) в виде номеров цилиндра, головки, сектора.
	\item Номер начального сектора LBA(4 байта) раздела.
	\item Количество секторов LBA(4 байта) раздела.
	\end{enumerate}
	
\item Затем идут сектора четырёх разделов, которые содержат структуры своих файловых систем.
\end{enumerate}

\subsection{Структура раздела с файловой системой FAT-32}

В записи раздела в таблице разделов главной загрузочной записи устанавливается тип файловой системы = 12($C_{16}$).

Раздел с файловой системой FAT-32 состоит из:

\begin{itemize}
\item зарезервированной области;
\item области таблиц FAT;
\item области данных.
\end{itemize}

\subsubsection{Зарезервированная область файловой системы FAT-32}

Зарезервированная область(Reserved Region) содержит следующую информацию:

\begin{itemize}
\item 0 сектор раздела -- Блок параметров BIOS(BPB -- BIOS Parameter Block);
\item 1 сектор раздела -- структура FSInfo;
\item 2-5 сектор раздела -- зарезервированные сектора;
\item 6 сектор раздела -- копия BPB;
\item 7 сектор раздела -- копия структуры FSInfo;
\item далее идут зарезервированные сектора, количество которых указывается в BPB.
\end{itemize}

\paragraph{Блок параметров BIOS}

Блок параметров BIOS(BPB) содержит следующую информацию:

\begin{enumerate}
\item Команда перехода JMP(3 байта) на загрузчик операционной системы. Может принимать 2 вида значений:

	\begin{enumerate}
	\item 1 байт равен 233($E9_{16}$), затем 1 байт короткого смещения, а после команда NOP(байт равен 144($90_{16}$)).
	\item 1 байт равен 235($EB_{16}$), затем 2 байта длинного смещения.
	\end{enumerate}
	
Длинное смещение используется, если загрузчик расположен в зарезервированных секторах.
\item Идентификатор версии операционной системы(8 байт) в виде строки.
\item Размер одного сектора(2 байта) в байтах.
\item Размер одного блока(1 байт) в секторах.
\item Количество секторов(2 байта) зарезервированной области.
\item Количество таблиц FAT(1 байт).
\item Количество записей корневого каталога(2 байта). Не используется в файловой системе FAT-32. Имеет значение равное 0.
\item Количество секторов в разделе(2 байта), если количество секторов меньше 65536, иначе равно 0.
\item Медиа дескриптор(1 байт). Хранит значение, определяющее тип диска.
\item Размер таблицы FAT(2 байта) в секторах. Не используется в файловой системе FAT-32. Имеет значение равное 0.
\item Количество секторов в дорожке(2 байта).
\item Количество головок диска(2 байта).
\item Количество скрытых секторов(4 байта) перед началом раздела.
\item Количество секторов в разделе(4 байта), если количество секторов больше или равно 65536.
\item Количество секторов на одну таблицу FAT(4 байта).
\item Номер активной таблицы FAT(2 байта). Побитовое значение:

	\begin{itemize}
	\item 0-3 бит -- номер активной FAT;
	\item 4-6 бит -- зарезервированы;
	\item 7 бит -- 0, если все копии FAT копируют активную FAT, иначе 1.
	\item 8-15 бит -- зарезервированы.
	\end{itemize}
	
\item Номер версии FAT(2 байта). Старший байт -- основной номер версии. Младший байт -- младшая версия. 
\item Номер первого блока корневого каталога(4 байта) в области данных. Обычно имеет значение равное 2.
\item Номер сектора активной структуры FSInfo(2 байта). 
\item Номер сектора копии BPB(2 байта).
\item Зарезервированные байты(12 байт).
\item Номер диска(1 байт).
\item Зарезервированный байт(1 байт).
\item Сигнатура(1 байт) наличия серийного номера или метки раздела. Имеет значение равное 41($29_{16}$), если значение номера или метки раздела не равно 0, иначе имеет значение равное 40($28_{16}$).
\item Серийный номер раздела(4 байта). Обычно имеет значение равное дате и времени при форматировании раздела.
\item Метка раздела(11 байт). Имеет значение равное значению, которое записано в особом файле родительского каталога с атрибутом метки раздела.
\item Строка идентификатора системы(8 байт). Обычно имеет значение \textquotedbl FAT32~~~\textquotedbl\
\item Зарезервированные байты(420 байт).
\item Сигнатура конца(2 байта). Имеет значение 43605($AA55_{16}$).
\end{enumerate}

\paragraph{Структура FSInfo}

Структура FSInfo содержит следующую информацию:

\begin{enumerate}
\item Сигнатура структуры FSInfo(4 байта). Имеет значение равное 1096897106($41615252_{16}$).
\item Зарезервированные байты(480 байт).
\item Сигнатура целостности структуры FSInfo(4 байта). Имеет значение равное 1631679090($61417272_{16}$).
\item Текущее количество свободных блоков(4 байта). Если имеет значение равное 4294967295($FFFFFFFF_{16}$), то это значение необходимо вычислить.
\item Номер первого свободного блока(4 байта). Указывает на номер блока, с которого необходимо искать свободный блок. Если имеет значение равное 4294967295($FFFFFFFF_{16}$), то необходимо искать с начала области данных(2 блок).
\item Зарезервированные байты(12 байт).
\item Сигнатура конца(4 байта). Имеет значение 2857697280($AA550000_{16}$).
\end{enumerate}

\subsubsection{Область таблиц FAT файловой системы FAT-32}

Область FAT(FAT Region) состоит из:

\begin{enumerate}
\item Таблицы FAT.
\item Копии таблицы FAT. Количество копий указывается в BPB.
\end{enumerate}

\paragraph{Таблица FAT файловой системы FAT-32}

Таблица FAT(File Allocation Table) -- таблица, представляющая собой последовательность элементов. Каждый элемент представляет собой состояние блока данных в области данных.

Размер таблицы FAT указывается в BPB.

Первые два элемента таблицы FAT являются зарезервированными и имеют значения:

\begin{enumerate}
\item Первый(0) элемент таблицы FAT имеет значение: старшие 3 байта имеют значение 1048575($0FFFFF_{16}$), а младший байт имеет значение равное медиа дескриптору из BPB.
\item Второй(1) элемент таблицы FAT имеет значение: старшие 2 бита являются флагами для некорректного завершения работы с файловой системой(старший бит -- ошибка при размонтировании диска, младший бит -- ошибка при чтении или записи).
\end{enumerate}

Каждый элемент таблицы FAT занимает 4 байта(32 бита), из который старшие 4 бита зарезервированы, и может принимать следующие значения:

\begin{enumerate}
\item 0 -- блок свободен.
\item От 2 до количества блоков в области данных -- блок занят. Значение в таблице -- номер следующего блока файла.
\item От количества блоков в области данных + 1 до 268435446($FFFFFF6_{16}$) -- зарезервированные значение и не могут быть использованы.
\item 268435447($FFFFFF7_{16}$) -- блок повреждён.
\item От 268435448($FFFFFF8_{16}$) до 268435454($FFFFFFE_{16}$) -- зарезервированные значение и не рекомендованы к использованию, но могут быть интерпретированы как то, что блок занят и является последним блоком файла.
\item 268435455($FFFFFFF_{16}$) -- блок занят и является последним блоком файла.
\end{enumerate}

\paragraph{Определение списка блоков файла по таблице FAT файловой системы FAT-32}

Зная номер первого блока файла, в таблице FAT находят элемент с номером первого блока файла и смотрят его значение, если блок не является конечным, то смотрят в таблице элемент с этим номером, и так пока не найдётся элемент таблицы со значением конца файла.

В результате получается односвязный список блоков искомого файла.

\subsubsection{Область данных файловой системы FAT-32}

Область данных(Data Region) -- множество блоков, в которых хранятся файлы и каталоги. Нумерация блоков области данных начинается с 2.

Каждый файл и каталог описывается 32-байтной записью. Для родительского каталога нет своей записи, номер первого блока родительского каталога хранится в BPB. В родительском каталоге есть файл с особым атрибутом метки раздела.

Каждый каталог, кроме корневого будет иметь 2 файла нулевого размера с названиями \textquotedbl.\textquotedbl и \textquotedbl..\textquotedbl, которые являются ссылками на текущий и родительский каталог соответственно. Эти файлы будут иметь значение атрибута равным 23($17_{16}$). Значение номера первого блока содержимого будет первым блоком текущего и родительского каталога соответственно.

\paragraph{Запись в каталоге файловой системы FAT-32}

Каждая 32-байтная запись в каталоге содержит следующую информацию:

\begin{enumerate}
\item Короткое имя файла или каталога(11 байт) -- названия файла в формате 8+3(8 байт на имя, 3 байта на расширение). Для файла с атрибутом метки раздела значение имени копируется в поле метки раздела в BPB. Допустимые символы в коротком имени:

	\begin{itemize}
	\item цифры;
	\item заглавные буквы латиницы;
	\item специальные символы: $\$~\%~'~-~\_~@~\sim~`~!~(~)~\{~\}~\wedge~\#~\&$.
	%\begin{verbatim}$ % ' - _ @ ~ ` ! ( ) { } ^ # &\end{verbatim} 
	\end{itemize}
	
Если первый символ короткого имени файла имеет значение 229($E5_{16}$), то это значит, что запись в каталоге помечена как удалённая.
\item Атрибуты файла(1 байт). Побитовые атрибуты:

	\begin{itemize}
	\item 0 бит -- только для чтения;
	\item 1 бит -- скрытый файл;
	\item 2 бит -- системный файл;
	\item 3 бит -- метка раздела;
	\item 4 бит -- директория;
	\item 5 бит -- изменённый(если файл или каталог был создан, переименован или модифицирован).
	\end{itemize}
	
\item Зарезервированный байт(1 байт).
\item Сотые доли секунды создания файла или каталога(1 байт). Имеет значение от 0 до 199(от 0 до 1.99 секунды с шагом 0.01 секунда).
\item Время создания файла или каталога(2 байта). Имеет значение в формате:

	\begin{itemize}
	\item 0-4 бит -- секунда делённые на 2;
	\item 5-10 бит -- минута;
	\item 11-15 -- час.
	\end{itemize}
	
\item Дата создания файла или каталога(2 байта). Имеет значение в формате:

	\begin{itemize}
	\item 0-4 бит -- день;
	\item 5-8 бит -- месяц;
	\item 9-15 -- год.
	\end{itemize}
	
\item Дата последнего доступа к файлу или каталогу(2 байта). Имеет такой же формат как и дата создания.
\item Старшие 2 байта номера первого блока содержимого файла или каталога.
\item Время последнего изменения файла или каталога(2 байта). Имеет такой же формат как и время создания.
\item Дата последнего изменения файла или каталога(2 байта). Имеет такой же формат как и дата создания.
\item Младшие 2 байта номера первого блока содержимого файла или каталога.
\item Размер файла в байтах(4 байта).
\end{enumerate}

\paragraph{Запись длинного имени файловой системы FAT-32}

Если для имени не хватает места в записи в каталоге, то прямо перед записью в каталоге создаётся запись длинного имени, которая содержит следующую информацию:

\begin{enumerate}
\item Номер записи в последовательности записей длинного имени(1 байт). Нумерация должна иметь 7 бит равным 1(маска 64($40_{16}$)) -- это позволяет определить, что это запись длинного имени.
\item Первые 5 символов Unicode записи(10 байт).
\item Атрибут(1 байт) равный 16($F_{16}$) -- запись длинного имени.
\item Тип длинной записи(1 байт). Для записи длинного имени имеет значение равное 0.
\item Контрольная сумма короткого имени(1 байт).
\item Следующие 6 символов Unicode записи(12 байт).
\item Зарезервированные байты(2 байта).
\item Последние 2 символа Unicode записи(4 байта).
\end{enumerate}

Допустимые символы для длинного имени:

\begin{itemize}
\item цифры;
\item буквы Unicode;
\item пробелы;
\item специальные символы: $\$~\%~'~-~\_~@~\sim~`~!~(~)~\{~\}~\wedge~\#~\&~+~,~;~=~[~]$.
\end{itemize}

Если имя не полностью помещается в запись длинного имени, то в конце имени ставится терминальный символ(0) и остальные символы ставятся 255($FF_{16}$).

Расположение записей длинного имени перед записью в каталоге идёт в обратном порядке, то есть сначала идёт последняя запись, затем предпоследняя, и так до первой.

\paragraph{Конвертирование длинного имени в короткое имя файловой системы FAT-32}

Алгоритм конвертации длинного имени в короткое:

\begin{enumerate}
\item Конвертирует длинное имя в верхний регистр.
\item Находим аналогичные символы в ASCII, если символа нет, то он заменяется на \_.
\item Убираются все пробелы.
\item Убираются все начальные точки.
\item Взять первые 8(или меньше, если символов основного имени меньше 8) символа основного имени, корректных для короткого имени.
\item Найти последнюю точку в название перед расширением.
\item Взять первые 3(или меньше, если символов расширения меньше 3) символа расширения, корректных для короткого имени.
\end{enumerate}

Если хотя бы 1 символ в длинном имени был заменён на \_, или длинное имя полностью помещается в короткое имя, или полученное короткое имя уже есть в текущем каталоге, то в конце имени файла добавляется числовой хвост.

Алгоритм добавления числового хвоста:

\begin{enumerate}
\item Находится число для числового хвоста -- первое доступное число, среди числовых хвостов.
\item В конце основного имени файла добавляется символ $\sim$, а после него число для числового хвоста, если длина основного имени не позволяется добавить числовой хвост, то конец основного имени заменяется на числовой хвост.
\end{enumerate}

Полученное имя является коротким именем конвертированным из длинного имени.

\subsubsection{Создание нового файла или каталога файловой системы FAT-32}

Порядок действий при создании нового файла или каталога:

\begin{enumerate}
\item Высчитывается сколько записей длинного имени понадобится, чтобы вместить имя файла.
\item Создаются записи длинного имени и в каждой устанавливается:

	\begin{itemize}
	\item номер записи в последовательности;
	\item символы имени;
	\item атрибут;
	\end{itemize}
	
\item Внутри содержимого текущего каталога создаётся новая запись в каталоге.
\item Если в блоке текущего каталога нет места, то выделяется новый блок, а в таблице FAT, в элементе с номером последнего блока текущего каталога, устанавливается номер выделенного блока.
\item В созданной записи устанавливается:

	\begin{itemize}
	\item короткое имя файла и его расширение;
	\item если создаётся каталог, то ставится атрибут каталога;
	\item устанавливается время и дата создания.
	\end{itemize}
	
\item Выделяется первый свободный блок, и его номер устанавливается в созданную запись.
\item В таблице FAT, в элемент с номером выделенного для содержимого файла или каталога блока, устанавливается значение конца файла.
\item Если создаётся каталог, то в новом каталоге создаются 2 файла нулевого размера с названиями \textquotedbl.\textquotedbl и \textquotedbl..\textquotedbl, в которых устанавливается соответствующий атрибут и значение первого блока содержимого.
\end{enumerate}

\subsubsection{Чтение файла или каталога файловой системы FAT-32}

Порядок действий при чтении файла или каталога:

\begin{enumerate}
\item Читается запись в каталоге искомого файла или каталога.
\item Из записи в каталоге узнаётся номер первого блока файла или каталога, а также его размер.
\item Находится нужный блок и читается.
\item В таблице FAT находится элемент с номером блока и читается его значение: если значение элемента не является значением конца файла, то находится номер следующего блока, читается новый блок и находится номер следующего. Процесс продолжается, пока значение элемента таблицы FAT не будет значением конца файла, этот блок будет последним для файла.
\end{enumerate}

\subsubsection{Запись в файл файловой системы FAT-32}

Порядок действий при записи в файл:

\begin{enumerate}
\item Читается запись в каталоге искомого файла.
\item Из записи в каталоге узнаётся номер первого блока файла.
\item Если позиция в которую нужно записывать находится не в этом блоке, то в таблице FAT находится элемент с номером блока и читается его значение: если позиция находится не в этом блоке, то находится номер следующего блока, и так до тех пор пока не найдётся блок с нужной позицией для записи.
\item Начиная с позиции для записи, записываются байты в файл.
\item Если блок закончился, то в таблице FAT находится элемент с номером текущего блока и читается его значение: если значение элемента не является значением конца файла, то находится номер следующего блока, и продолжается запись. Иначе, если значение элемента в таблице имеет значение конца файла, то выделяется новый блок, а в элемент таблицы FAT записывается номер нового блока, и начинается запись в новый блок. Процесс продолжается пока не закончатся свободные блоки, или не закончится запись.
\end{enumerate}

\subsection{Структура раздела с файловой системой ISO-9660(CDFS)}

В записи раздела в таблице разделов главной загрузочной записи устанавливается тип файловой системы = 148($94_{16}$).

Раздел с файловой системой ISO-9660(CDFS) состоит из:

\begin{itemize}
\item системной области;
\item области дескрипторов тома;
\item двух таблиц путей;
\item области данных.
\end{itemize}

\subsubsection{Формат данных Little-Endian и Big-Endian файловой системы ISO-9660(CDFS)}

Данные в файловой системе ISO-9660(CDFS) могут хранится в 3 видах:

\begin{enumerate}
\item Little-Endian -- порядок байт от младшего к старшему(справа налево).
\item Big-Endian -- порядок байт от старшего к младшему(слева направо).
\item Both-Endian -- сначала идёт значение в Little-Endian, а после в Big-Endian.
\end{enumerate}

\subsubsection{Системная область файловой системы ISO-9660(CDFS)}

Системная область(System Area) занимает первые 64 сектора раздела(32КБ), содержимое которого зарезервировано.

\subsubsection{Область дескрипторов тома файловой системы ISO-9660(CDFS)}

Область дескрипторов тома(Volume Descriptors) состоит из:

\begin{itemize}
\item загрузочной записи дескриптора тома;
\item основного дескриптора тома;
\item ограничителя набора дескрипторов тома.
\end{itemize}

Каждый дескриптор тома занимает 4 сектора(2048 байт).

\paragraph{Загрузочная запись дескриптора тома файловой системы ISO-9660(CDFS)}

Загрузочная запись(The Boot Record) содержит следующую информацию:

\begin{enumerate}
\item Тип дескриптора тома(1 байт). Имеет значение равное 0.
\item Идентификатор дескриптора тома(5 байт). Имеет значение равное \textquotedbl CD001\textquotedbl.
\item Версия дескриптора тома(1 байт). Имеет значение равное 1.
\item Идентификатор загрузочной системы(32 байта). Содержит идентификатор системы, которая может запустить систему из загрузчика. Если не используется, то заполняется пробелами.
\item Идентификатор загрузчика(32 байта).
\item Содержимое загрузчика системы(1977 байт). 
\end{enumerate}

\paragraph{Основной дескриптор тома файловой системы ISO-9660(CDFS)}

Основной дескриптор тома(The Primary Volume Descriptor) содержит следующую информацию:

\begin{enumerate}
\item Тип дескриптора тома(1 байт). Имеет значение равное 1.
\item Идентификатор дескриптора тома(5 байт). Имеет значение равное \textquotedbl CD001\textquotedbl.
\item Версия дескриптора тома(1 байт). Имеет значение равное 1.
\item Зарезервированный байт(1 байт).
\item Идентификатор системы(32 байта). Содержит идентификатор системы, которая может взаимодействовать с системной областью. Если не используется, то заполняется пробелами.
\item Идентификатор тома(32 байта). Название раздела.
\item Зарезервированные байты(8 байт).
\item Количество блоков раздела(8 байт). Имеет значение в формате Both-Endian.
\item Зарезервированные байты(32 байта).
\item Размер набора томов. Имеет значение равное 1. Имеет значение в формате Both-Endian.
\item Порядковый номер тома. Имеет значение равное 1. Имеет значение в формате Both-Endian.
\item Размер одного блока в байтах(4 байта). Имеет значение в формате Both-Endian.
\item Размер таблицы путей в байтах(8 байт). Имеет значение в формате Both-Endian.
\item Номер блока, в котором находится таблица путей(4 байта) в формате Little-Endian. Имеет значение в формате Little-Endian.
\item Номер блока, в котором находится копия таблица путей(4 байта) в формате Little-Endian. Имеет значение в формате Little-Endian. Если копии нет, то имеет значение равное 0.
\item Номер блока, в котором находится таблица путей(4 байта) в формате Big-Endian. Имеет значение в формате Big-Endian.
\item Номер блока, в котором находится копия таблица путей(4 байта) в формате Big-Endian. Имеет значение в формате Big-Endian. Если копии нет, то имеет значение равное 0.
\item Запись в каталоге для корневого каталога(34 байта). Имеет длину имени равную 1. Байт имени имеет значение равное 0.
\item Идентификатор набора томов(128 байт). Имеет значение равное идентификатору тома.
\item Идентификатор издателя(128 байт). Если не используется, то заполняется пробелами.
\item Идентификатор поставщика данных(128 байт). Если не используется, то заполняется пробелами.
\item Идентификатор расширения(128 байт). Если не используется, то заполняется пробелами.
\item Идентификатор файла с авторскими правами(37 байт). Если не используется, то заполняется пробелами.
\item Идентификатор файла с аннотацией(37 байт). Если не используется, то заполняется пробелами.
\item Идентификатор библиографического файла(37 байт). Если не используется, то заполняется пробелами.
\item Дата и время создания раздела(17 байт).
\item Дата и время модификации раздела(17 байт).
\item Дата и время истечения срока раздела(17 байт). После этой даты и времени этот раздел нельзя будет использовать. Если даты нет, то имеет значение равное 0.
\item Дата и время активации раздела(17 байт). После этой даты и времени этот раздел можно будет использовать. Если даты нет, то имеет значение равное 0.
\item Версия структуры файлов(1 байт). Имеет значение равное 1.
\item Зарезервированный байт(1 байт).
\item Расширение стандарта ISO-9660(512 байт).
\item Зарезервированные байты(653 байта).
\end{enumerate}

\paragraph{Ограничитель набора дескрипторов тома файловой системы ISO-9660(CDFS)}

Ограничитель набора дескрипторов тома(Volume Descriptor Set Terminator) содержит следующую информацию:

\begin{enumerate}
\item Тип дескриптора тома(1 байт). Имеет значение равное 255($FF_{16}$).
\item Идентификатор дескриптора тома(5 байт). Имеет значение равное \textquotedbl CD001\textquotedbl.
\item Версия дескриптора тома(1 байт). Имеет значение равное 1.
\end{enumerate}

\subsubsection{Формат даты и времени в основном дескрипторе тома файловой системы ISO-9660(CDFS)}

Дата и время в основном дескрипторе тома содержит следующую информацию:

\begin{enumerate}
	\item Год(4 байта). От 1 до 9999. Имеет значение в виде строки.
	\item Месяц(2 байта). От 1 до 12. Имеет значение в виде строки.
	\item День(2 байта). От 1 до 31. Имеет значение в виде строки.
	\item Час(2 байта). От 0 до 23. Имеет значение в виде строки.
	\item Минута(2 байта). От 0 до 59. Имеет значение в виде строки.
	\item Секунда(2 байта). От 0 до 59. Имеет значение в виде строки.
	\item Сотая секунды(2 байта). От 0 до 99. Имеет значение в виде строки.
	\item Смещение часового пояса(1 байт). Исчисляется в 15 минутных интервалах от -48 до 52(0 = -48 = -12 часов, 100 = 52 = +13 часов). Имеет значение в виде числа.
\end{enumerate}

\subsubsection{Таблица путей файловой системы ISO-9660(CDFS)}

Таблица путей(The Path Table) хранится в 2 форматах: Little-Endian и Big-Endian. Разница между ними в том, что значения внутри таблиц хранятся в соответствующем формате. Местоположение таблиц путей указывается в основном дескрипторе тома.

Сама таблица путей -- последовательность записей таблицы путей. Каждая запись таблицы путей содержит следующую информацию:

\begin{enumerate}
\item Длина имени каталога(1 байт). Для родительского каталога имеет значение равное 1.
\item Длина записи расширенных атрибутов(1 байт). Если не используется, то имеет значение равное 0.
\item Номер блока(4 байта) содержимого каталога.
\item Индекс родительского каталога в таблице путей(2 байта).
\item Имя каталога(длина хранится выше). Для родительского каталога имеет значение равное 0. 
\item Если запись таблицы путей имеет нечётную длину, то в конце записи таблицы путей добавляется байт со значением равным 0.
\end{enumerate}

\subsubsection{Область данных файловой системы ISO-9660(CDFS)}

Область данных(Data Area) -- множество блоков, в которых хранятся файлы и каталоги. Нумерация блоков области данных начинает с 16. Содержимое каталогов описывается последовательностью записей каталога. Для родительского каталога нет своей записи, номер записи таблицы путей корневого каталога, из которой можно узнать его номер блока, хранится в основном дескрипторе тома.

Каждый каталог, кроме корневого будет иметь 2 файла нулевого размера с названиями \textquotedbl\textquotedbl(\textquotedbl.\textquotedbl) и \textquotedbl\textbackslash1\textquotedbl(\textquotedbl..\textquotedbl), которые являются ссылками на текущий и родительский каталог соответственно. Значение номера первого блока содержимого будет первым блоком текущего и родительского каталога соответственно.

Всё содержимое файла или каталога записывается в виде экстента(непрерывной последовательности блоков).

Каждая запись каталога содержит следующую информацию:

\begin{enumerate}
\item Длина записи каталога(1 байт).
\item Длина записи расширенных атрибутов(1 байт). Если не используется, то имеет значение равное 0.
\item Номер блока с содержимым файла или каталога(8 байт). Имеет значение в формате Both-Endian.
\item Размер содержимого в байтах(8 байт). Имеет значение в формате Both-Endian.
\item Дата и время записи(7 байт).
\item Атрибуты файла(1 байт). Побитовые атрибуты:

	\begin{itemize}
	\item 0 бит -- скрытый файл;
	\item 1 бит -- директория;
	\item 2 бит -- ассоциативный файл;
	\item 3 бит -- флаг наличия формата файла в записи расширенных атрибутов;
	\item 4 бит -- флаг наличия прав доступа в записи расширенных атрибутов;
	\item 5 и 6 бит -- зарезервированы;
	\item 7 бит -- не последняя запись каталога для файла(если файл больше 4 ГБ).
	\end{itemize}
	
\item Размер файлового юнита для пробела между записями(1 байт). Обычно имеет значение равное 0. 
\item Разбер промежутка между блоками(1 байт). Обычно имеет значение равное 0. 
\item Порядковый номер тома(4 байта). Обычно имеет значение равное 1. Имеет значение в формате Both-Endian.
\item Длина имена файла в байтах(1 байт).
\item Имя файла(длина хранится выше). Разрешённые символы:

	\begin{itemize}
	\item заглавные буквы латиницы;
	\item специальный символ \_.
	\end{itemize}
	
В конце имени добавляется \textquotedbl;1\textquotedbl
\item Если запись каталога имеет нечётную длину, то в конце записи каталога добавляется байт со значением равным 0.
\item Байты для системного использования расширениями стандарта ISO-9660.
\end{enumerate}

\subsubsection{Формат даты и времени в записи каталога файловой системы ISO-9660(CDFS)}

Дата и время в записи каталога(все значения хранятся в виде чисел) содержит следующую информацию:

\begin{enumerate}
\item Год(1 байт). Значение года после 1900.
\item Месяц(1 байт). От 1 до 12.
\item День(1 байт). От 1 до 31.
\item Час(1 байт). От 0 до 23.
\item Минута(1 байт). От 0 до 59.
\item Секунда(1 байт). От 0 до 59.
\item Смещение часового пояса(1 байт). Исчисляется в 15 минутных интервалах от -48 до 52.
\end{enumerate}

\subsubsection{Создание нового файла или каталога файловой системы ISO-9660(CDFS)}

Файловая система ISO-9660(CDFS) является файловой системой только для чтения, из-за чего создание нового файла или каталога невозможно.

\subsubsection{Чтение файла или каталога файловой системы ISO-9660(CDFS)}

Порядок действий при чтении файла или каталога:

\begin{enumerate}
\item Читается запись каталога искомого файла или каталога.
\item Из записи каталога узнаётся номер первого блока файла или каталога, а также его размер.
\item Находится нужный блок и читается.
\item Если блок закончился, то читается следующий блок до тех пор, пока файл или каталог не закончится.
\end{enumerate}

\subsubsection{Запись в файл файловой системы ISO-9660(CDFS)}

Файловая система ISO-9660(CDFS) является файловой системой только для чтения, из-за чего запись в файл невозможна.

\subsection{Общая архитектура системы}

Система включает в себя следующие модули:

\begin{enumerate}
\item Модуль командного интерпретатора. Этот модуль отвечает за командный интерпретатор.
\item Модуль интерфейса виртуальной файловой системы. Этот модуль отвечает за интерфейс виртуальной файловой системы.
\item Модуль главной загрузочной записи. Этот модуль отвечает за загрузку файловой системы из определённого раздела диска.
\item Модуль файловой системы FAT-32. Этот модуль отвечает за реализацию интерфейса файловой системы для файловой системы FAT-32.
\item Модуль файловой системы ISO-9660(CDFS). Этот модуль отвечает за реализацию интерфейса файловой системы для файловой системы ISO-9660(CDFS).
\item Модуль виртуальной файловой системы. Этот модуль отвечает за класс абстрактной виртуальной файловой систему и класс абстрактного файла.
\item Модуль работы с блоками диска. Этот модуль отвечает за чтение блоков диска, запись блоков на диск.
\item Модуль драйвера ATA. Этот модуль отвечает за общение с интерфейсом ATA: чтение секторов диска, запись секторов на диск.
\end{enumerate}

\subsection{Диаграмма компонентов}

На рисунке \ref{ComponentDiagram:image} изображена диаграмма компонентов для проектируемой
системы. Она включает модули:

\begin{itemize}
\item модуль командного интерпретатора;
\item модуль интерфейса виртуальной файловой системы;
\item модуль главной загрузочной записи;
\item модуль файловой системы FAT-32;
\item модуль файловой системы ISO-9660(CDFS);
\item модуль виртуальной файловой системы;
\item модуль работы с блоками диска;
\item модуль драйвера ATA.
\end{itemize}

\begin{figure}[H]
\center{\includegraphics[width=1\linewidth]{ComponentDiagram}}
\caption{Диаграмма компонентов}
\label{ComponentDiagram:image}
\end{figure}

\subsection{Модуль командного интерпретатора}

Командный интерпретатор позволяет принимать набранную с клавиатуры строку как команду с аргументами, что даёт возможность интерактивного взаимодействия с файловой системой.

Для реализации командного интерпретатора необходимо:

\begin{itemize}
\item создать обработчик нажатия клавиш с клавиатуры;
\item создать парсер команд и их аргументов;
\item создать обработчики команд с их аргументами.
\end{itemize}

Модуль командного интерпретатора включает следующие функции:

\begin{enumerate}
\item \texttt{key-down-handler} -- обработчик нажатия клавиш с клавиатуры. Аргументы: \texttt{scan}(скан код нажатой клавиши на клавиатуре). Возвращает \texttt{nil}.
\item \texttt{process-command} -- обработка команды. Аргументы: \texttt{com}(список введённых символов с клавиатуры). Возвращает \texttt{nil}.
\item \texttt{parse-commands} -- парсер команд. Аргументов нет. Возвращает парсер команд.
\item \texttt{parse-command} -- парсер определённой команды и её атрибутов. Аргументы: \texttt{com}(строка названия команды), \texttt{fun}(функция, которую выполняет команда). Возвращает парсер команды и её атрибутов.
\end{enumerate}

\subsection{Модуль интерфейса виртуальной файловой системы}

Пользователю не нужно знать, какая файловая система находится на разделе диска, поэтому можно выделить функции интерфейса виртуальной файловой системы, а виртуальная файловая система будет вызывать методы той файловой системы, которая находится на диске, для этого нужна функция загрузки файловой системы конкретного типа, которая создаст экземпляр класса файловой системы определённого типа и сохранит его в глобальной переменной \texttt{file-system}.

Каждая функция виртуальной файловой системы -- вызов соответствующего метода из \texttt{file-system}.

Функция открытия файла возвращает экземпляр класса файла конкретной файловой системы, который можно передать в качестве аргументы в функции работы с файлом. Реализация этого экземпляра зависит от конкретной файловой системы.

Пути в стиле UNIX(/dir/file.txt).

Структуру каталогов удобно хранить в естественном виде дерева из файлов и каталогов. Это будет файловый объект, хранящий необходимую информацию для того, чтобы писать в этот файл или каталог, у которого есть поле \texttt{dir}, которое, если это файл, то равно \texttt{nil}, если это нераскрытый каталог, то равно \texttt{t}, если это загруженный каталог, то это список файловых объектов. Каталоги раскрываются по мере чтения. Функция \texttt{is-directory} возвращает \texttt{t} если файл -- каталог. Глобальный переменная \texttt{root-directory} будет хранить файловый объект корневого каталога с деревом файлов и каталогов. Глобальная переменная \texttt{working-directory} является ссылкой на часть \texttt{root-directory}, хранящая файловый объект текущего каталога. Для раскрытия каталога используется метод виртуальной файловой системы \texttt{load-dir}.

Пример дерева файлов и каталогов:

\begin{figure}[H]
	\begin{lstlisting}[language=Lisp]
    Дерево корневого каталога:
    <Файловый объект
    (dir . (
            ("File.txt" . <Файловый объект ; обычный файл
              (dir . nil)>)
            ("UnopenedDirectory" . <Файловый объект ; нераскрытый каталог
              (dir . t)>)
            ("OpenedDirectory" . <Файловый объект ; загруженный каталог
              (dir . (
                      ("File2.txt" . <Файловый объект
                        (dir . nil)>)
                     ))>))
           )>
	\end{lstlisting}
\end{figure}

Для того чтобы из строки пути получать дерево каталога требуется функция \texttt{load-path}, которая будет при необходимости раскрывать каталоги по пути.

Модуль виртуальной файловой системы включает следующие переменные:

\begin{enumerate}
	\item \texttt{file-system} -- экземпляр класса конкретной файловой системы.
	\item \texttt{root-directory} -- дерево файлов и каталогов файловой системы.
	\item \texttt{working-directory} -- дерево текущего каталога.
\end{enumerate}

Модуль виртуальной файловой системы включает следующие функции:

\begin{enumerate}
	\item \texttt{cur-dir} -- получить путь текущего каталога. Аргументов нет. Возвращает полный путь текущего каталога в виде строки.
	\item \texttt{list-dir} -- просмотреть содержимое каталога. Аргументы: \texttt{path}(путь к каталогу, содержимое которого необходимо вернуть). Возвращает список содержимого каталога, каждый элемент списка -- пара, состоящая из символа типа(файл, каталог) и строки названия.
	\item \texttt{create-dir} -- создать каталог. Добавить каталог в дерево файлов и каталогов. Аргументы: \texttt{path}(путь к каталогу, в котором создать каталог), \texttt{name}(название каталога). Возвращает \texttt{nil}.
	\item \texttt{remove-dir} -- удалить каталог, если он пуст. Удалить каталог из дерева файлов и каталогов. Аргументы: \texttt{path}(путь к каталогу, который нужно удалить). Возвращет \texttt{nil}.
	\item \texttt{change-dir} -- переместиться в каталог. При необходимости раскрыть каталог в дереве файлов и каталогов. Аргументы: \texttt{path}(путь к каталогу, в который необходимо переместиться). Возвращает \texttt{nil}.
	\item \texttt{create-file} -- создать и открыть файл. Добавить файл в дерево файлов и каталогов. Аргументы: \texttt{path}(путь к каталогу, в котором необходимо создать файл), \texttt{name}(название файла). Возвращает объект файла.
	\item \texttt{remove-file} -- удалить файл. Удалить файл в дереве файлов и каталогов. Аргументы: \texttt{path}(путь к файлу, который нужно удалить). Возвращет \texttt{nil}.
	\item \texttt{open-file} -- открыть файл как поток для чтения и записи. При необходимости раскрыть каталоги по пути к файлу в дереве файлов и каталогов. Аргументы: \texttt{path}(путь к файлу, который нужно открыть). Возвращает объект файла.
	\item \texttt{close-file} -- закрыть файл. Аргументы: \texttt{file-object}(объект файла, который необходимо закрыть). Возвращает \texttt{nil}.
	\item \texttt{read-file} -- прочитать определённое количество байт из потока файла, и переместить позицию внутри потока на то же количество. Аргументы: \texttt{file-object}(объект файла, из которого читать байты), \texttt{size}(количество байт, которое необходимо прочитать). Возвращает массив прочитанных байт.
	\item \texttt{write-file} -- записать массив байт в поток файла, и переместить позицию внутри потока на то же количество. Аргументы: \texttt{file-object}(объект файла, в который необходимо записать массив байт), \texttt{buffer}(массив байт, который необходимо записать). Возвращает количество байт, которые было записано.
	\item \texttt{tell-file} -- получить позицию внутри потока файла. Аргументы: \texttt{file-object}(объект файла, позицию в котором необходимо вернуть).
	\item \texttt{seek-file} -- переместить позицию внутри потока файла. Аргументы: \texttt{file-object}(объект файла, позицию в котором необходимо переместить), \texttt{offset}(насколько сместить позицию), \texttt{origin}(откуда смещать позицию). Возвращает значение новой позиции в потоке.
	\item \texttt{is-directory} -- проверить, является ли файл каталогом. Аргументы: \texttt{path}(путь к файлу, который надо проверить). Возвращает \texttt{t}, если файл является каталогом, и \texttt{nil}, если нет.
	\item \texttt{rename} -- переименовать файл или каталог. Аргументы: \texttt{path}(путь к файлу или каталогу, который необходимо переименовать), \texttt{name}(новое имя файла или каталога). Возвращает \texttt{nil}.
	\item \texttt{free-space} -- получить объём свободного места в текущем разделе диска. Аргументов нет. Возвращает размер в байтах.
	\item \texttt{set-attr} -- изменить атрибуты файла или каталога. Аргументы: \texttt{path}(путь к файлу или каталогу, чьи атрибуты необходимо поменять), \texttt{new-attributes}(новые атрибуты файла или каталога). Возвращает \texttt{nil}.
	\item \texttt{fstat} -- получить метаданные файла или каталога. Аргументы: \texttt{path}(путь к файлу или каталогу, чьи метаданные надо получить). Возвращает структуру \texttt{stat}. Структура \texttt{stat} состоит из:
	
	\begin{enumerate}
		\item \texttt{name} -- название файла;
		\item \texttt{size} -- размер файла в байтах. Для каталогов равно 0;
		\item \texttt{create-date} -- дата создания файла, если нет в конкретной файловой системе, то равно \texttt{nil};
		\item \texttt{create-time} -- время создания файла, если нет в конкретной файловой системе, то равно \texttt{nil};
		\item \texttt{modify-date} -- дата последнего изменения файла, если нет в конкретной файловой системе, то равно \texttt{nil};
		\item \texttt{modify-time} -- время последнего изменения файла, если нет в конкретной файловой системе, то равно \texttt{nil};
		\item \texttt{access-date} -- дата последнего доступа файла, если нет в конкретной файловой системе, то равно \texttt{nil};
		\item \texttt{access-time} -- время последнего изменения файла, если нет в конкретной файловой системе, то равно \texttt{nil};
		\item \texttt{dir} -- флаг, показывающий, является ли файл каталогом;
		\item другие атрибуты, которые зависят от конкретной файловой системы.
	\end{enumerate}
	
	\item \texttt{load-path} -- получить дерево каталога по абсолютному или относительному пути. При необходимости раскрыть каталоги по пути к каталогу в дереве файлов и каталогов. Аргументы: \texttt{path}(путь к каталогу). Возвращает дерево каталога по пути.
		
\end{enumerate}

\subsection{Модуль главной загрузочной записи}

Для того, чтобы загрузить файловую систему из определённого раздела, необходимо найти нужную запись раздела из таблицы разделов, получить номер начального сектора и количество секторов и вызвать функцию загрузки файловой системы определённого типа из модуля виртуальной файловой системы.

Модуль главной загрузочной записи включает следующую функцию:

\texttt{load-partition} -- загрузить файловую систему из определённого раздела диска. Аргументы: \texttt{disk-num}(номер диска, с которого загрузить файловую систему), \texttt{part-num}(номер раздела, с которого загрузить файловую систему). Возвращает \texttt{nil}.

\subsection{Модуль виртуальной файловой системы}

Виртуальная файловая система -- два класса, методы которого являются функциями интерфейса виртуальной файловой системы.

Поскольку функции \texttt{open-file}, \texttt{close-file}, \texttt{read-file}, \texttt{write-file}, \texttt{seek-file}, \texttt{tell-file} работают с файлами, то функция \texttt{open-file} будет возвращать файловый объект, а остальные будут работать с этим файловым объектом. Из-за этого стоит выделить класс файлового объекта.

Поскольку чтение и запись файла по блокам одинакова для всех файловых систем, то можно реализовать методы \texttt{read-file}, \texttt{write-file}, \texttt{seek-file}, \texttt{tell-file}, а в реализациях конкретных файловых систем вызывать методы родительского класса и менять метаданные файлов и каталогов.

Поле \texttt{blocks} в классе файла -- список блоков файла.

Поле \texttt{position} в классе файла -- пара из номера текущего блока и позиции в нём.

Если во время чтения позиция вышла за текущий блок, то проверяется есть ли у этого файла следующий блок, если блока нет, то вызывается ошибка, иначе номер следующего блока записывается в поле \texttt{position}, позиция в текущем блоке в поле \texttt{position} обнуляется.

При записи в файл может понадобится добавление нового блока к файлу, для этого стоит добавить метод \texttt{new-block}, реализация которого будет в конкретной файловой системе.

Для остальных функций интерфейса виртуальной файловой выделяется класс файловой системы, реализация функций будет в конкретной файловой системе.

Для того чтобы раскрывать каталога в дереве файлов и каталогов добавляется метод \texttt{load-dir}, реализация которого будет в конкретной файловой системе.

Классы виртуальной файловой системы:

\begin{enumerate}
	\item Класс абстрактной файловой системы. Этот класс отвечает за создание, удаление, переименование, открытие файлов и каталогов.
	\item Класс абстрактного файла. Этот класс отвечает за чтение из файла и запись в него.
\end{enumerate}

Каждая файловая система -- два класса, являющиеся наследниками классов виртуальной файловой системы, которые реализуют функции интерфейса виртуальной файловой системы.

Модуль виртуальной файловой системы включает следующие классы:

\begin{enumerate}
	\item \texttt{FileSystem} -- класс абстрактной файловой системы. Методы класса \texttt{FileSystem}:
	
	\begin{enumerate}
		\item \texttt{fs-init} -- загрузка файловой системы. Аргументы: \texttt{disk}(номер диска, с которого загружают), \texttt{start-sector}(номер сектора начала раздела), \texttt{sector-count}(количество секторов раздела). Возвращает \texttt{nil}.
		\item \texttt{create-dir} -- создать каталог. Аргументы: \texttt{dir}(каталог, в котором создать каталог), \texttt{name}(название каталога). Возвращает \texttt{nil}.
		\item \texttt{remove-dir} -- удалить каталог, если он пуст. Аргументы: \texttt{parent-dir}(каталог, из которого нужно удалить), \texttt{dir}(каталог, который нужно удалить). Возвращет \texttt{nil}.
		\item \texttt{load-dir} -- раскрыть и сохранить содержимое каталога. Аргументы: \texttt{dir}(каталог, чьё содержимое нужно загрузить). Возвращает \texttt{nil}.
		\item \texttt{create-file} -- создать и открыть файл. Аргументы: \texttt{dir}(каталог, в котором необходимо создать файл), \texttt{name}(название файла). Возвращает объект файла.
		\item \texttt{remove-file} -- удалить файл. Аргументы: \texttt{dir}(каталог, из которого нужно удалить), \texttt{entry}(файл, который нужно удалить). Возвращет \texttt{nil}.
		\item \texttt{open-file} -- открыть файл как поток для чтения и записи. Аргументы: \texttt{file}(файл, который нужно открыть). Возвращает объект файла.
		\item \texttt{rename} -- переименовать файл или каталог. Аргументы: \texttt{entry}(файл или каталог, который нужно переименовать), \texttt{new-name}(новое имя файла или каталога). Возвращает \texttt{nil}.
		\item \texttt{free-space} -- получить объём свободного места в текущем разделе диска. Аргументов нет. Возвращает размер в байтах.
		\item \texttt{set-attr} -- изменить атрибуты файла или каталога. Аргументы: \texttt{entry}(файл или каталог, чьи атрибуты необходимо поменять), \texttt{new-attributes}(новые атрибуты файла или каталога). Возвращает \texttt{nil}.
		\item \texttt{fstat} -- получить метаданные файла или каталога. Аргументы: \texttt{entry}(файл или каталог, чьи метаданные надо получить). Возвращает структуру \texttt{stat}. Структура \texttt{stat} состоит из:
		
		\begin{enumerate}
			\item \texttt{name} -- название файла;
			\item \texttt{size} -- размер файла в байтах. Для каталогов равно 0;
			\item \texttt{create-date} -- дата создания файла, если нет в конкретной файловой системе, то равно \texttt{nil};
			\item \texttt{create-time} -- время создания файла, если нет в конкретной файловой системе, то равно \texttt{nil};
			\item \texttt{modify-date} -- дата последнего изменения файла, если нет в конкретной файловой системе, то равно \texttt{nil};
			\item \texttt{modify-time} -- время последнего изменения файла, если нет в конкретной файловой системе, то равно \texttt{nil};
			\item \texttt{access-date} -- дата последнего доступа файла, если нет в конкретной файловой системе, то равно \texttt{nil};
			\item \texttt{access-date} -- время последнего изменения файла, если нет в конкретной файловой системе, то равно \texttt{nil};
			\item \texttt{isdir} -- флаг, показывающий, является ли файл каталогом;
			\item другие атрибуты, которые зависят от конкретной файловой системы.
		\end{enumerate}
		
		\item \texttt{new-block} -- добавить новый блок к файлу. Аргументы: \texttt{file-object}(объект файла, к которому необходимо добавить новый блок). Возвращает \texttt{nil}.
		
	\end{enumerate}
	
	\item \texttt{File} -- класс абстрактного файла. Поля класса \texttt{File}:
	
	\begin{enumerate}
		\item \texttt{name} -- имя файла.
		\item \texttt{size} -- размер файла в байтах.
		\item \texttt{position} -- позиция в потоке чтения/записи файла(пара из номера текущего блока и позиции в нём).
		\item \texttt{blocks} -- список номеров блока файла.
	\end{enumerate}
	
	Методы класса \texttt{File}:
	
	\begin{enumerate}
		\item \texttt{close-file} -- закрыть файл. Аргументы: \texttt{file-object}(объект файла, который необходимо закрыть). Возвращает \texttt{nil}.
		\item \texttt{read-file} -- прочитать определённое количество байт из потока файла, и переместить позицию внутри потока на то же количество. Аргументы: \texttt{file-object}(объект файла, из которого читать байты), \texttt{size}(количество байт, которое необходимо прочитать). Возвращает массив прочитанных байт.
		\item \texttt{write-file} -- записать массив байт в поток файла, и переместить позицию внутри потока на то же количество. При необходимости добавить новый блок к файлу. Аргументы: \texttt{file-object}(объект файла, в который необходимо записать массив байт), \texttt{buffer}(массив байт, который необходимо записать). Возвращает количество байт, которые было записано.
		\item \texttt{tell-file} -- получить позицию внутри потока файла. Аргументы: \texttt{file-object}(объект файла, позицию в котором необходимо вернуть).
		\item \texttt{seek-file} -- переместить позицию внутри потока файла. Аргументы: \texttt{file-object}(объект файла, позицию в котором необходимо переместить), \texttt{offset}(насколько сместить позицию), \texttt{origin}(откуда смещать позицию). Возвращает значение новой позиции в потоке.
	\end{enumerate}
	
\end{enumerate}

\subsection{Модуль файловой системы FAT-32}

Описание файловой системы FAT-32 приведено выше.

При загрузке файловой системы необходимо прочитать BPB, FSInfo, FAT.

Из BPB необходимо достать следующую информацию:

\begin{itemize}
\item размер блока в байтах;
\item количество секторов зарезервированной области;
\item размер таблицы FAT;
\item количество секторов на одну таблицу FAT;
\item номер активной таблицы FAT и режим копирования;
\item номер первого блока корневого каталога;
\item номер сектора активной структуры FSInfo.
\end{itemize}

Из FSInfo необходимо достать следующую информацию:

\begin{itemize}
\item количество свободных блоков(вычислить при необходимости);
\item номер, с которого необходимо начать поиск свободного блока(вычислить при необходимости).
\end{itemize}

Вычисление количества свободных блоков и номера, с которого необходимо начать поиск свободного блока происходит в момент выделения нового блока.

Для удобства можно хранить всю таблицу FAT, её удобно хранить как хеш-объект с ключом - номер первого блока в цепочке и значением -- список оставшихся блоков. При получении цепочки таблицы FAT -- возвращать цепочку из хеш-объекта, а при записи -- менять цепочку в хеш-объекте и записывать в таблицу FAT на диске.

При вызове функции открытия файла, создаётся экземпляр класса \texttt{FAT32File}. Список блоков файла сохраняется в поле \texttt{blocks}. Во время чтения или записи, позиция внутри текущего блока сохраняется в поле \texttt{position}, номер текущего блока сохраняется в поле \texttt{position}.

Модуль файловой системы FAT-32 включает следующие классы:

\begin{enumerate}
\item \texttt{FAT32FileSystem} -- класс файловой системы FAT-32. Наследник класса \texttt{FileSystem}. Поля класса \texttt{FAT32FileSystem}:

	\begin{enumerate}
	\item \texttt{fat-table} -- массив таблицы FAT.
	\item \texttt{fat-start-sector} -- номер сектора таблицы FAT.
	\item \texttt{fat-copy-start-sector} -- номер сектора копии таблицы FAT.
	\item \texttt{fat-sectors} -- количество секторов таблицы FAT.
	\item \texttt{free-blocks-count} -- количество свободных блоков.
	\item \texttt{free-block-num} -- номер, с которого необходимо начать поиск свободного блока.
	\item \texttt{root-dir-block} -- номер блока корневого каталога.
	\end{enumerate}

Методы класса \texttt{FAT32FileSystem}:

	\begin{enumerate}
	\item \texttt{fs-init} -- загрузка файловой системы. Чтение BPB, FSInfo, таблицы FAT. Аргументы: \texttt{disk}(номер диска, с которого загружают), \texttt{start-sector}(номер сектора начала раздела), \texttt{sector-count}(количество секторов раздела). Возвращает \texttt{nil}.
	\item \texttt{create-dir} -- создать каталог. Создать записи длинного имени, создать запись в каталоге, выделить блок для каталога. Аргументы: \texttt{dir}(каталог, в котором создать каталог), \texttt{name}(название каталога). Возвращает \texttt{nil}.
	\item \texttt{remove-dir} -- удалить каталог, если он пуст. Аргументы: \texttt{parent-dir}(каталог, из которого нужно удалить), \texttt{dir}(каталог, который нужно удалить). Возвращет \texttt{nil}.
	\item \texttt{load-dir} -- раскрыть и сохранить содержимое каталога. Аргументы: \texttt{dir}(каталог, чьё содержимое нужно загрузить). Возвращает \texttt{nil}.
	\item \texttt{create-file} -- создать файл. Создать записи длинного имени, создать запись в каталоге, выделить блок для файла. Аргументы: \texttt{dir}(каталог, в котором необходимо создать файл), \texttt{name}(название файла). Возвращает \texttt{nil}.
	\item \texttt{remove-file} -- удалить файл. Заменить первые символы во всех записях длинного имени, а также в записи каталога, на 229($E5_{16}$). Аргументы: \texttt{dir}(каталог, из которого нужно удалить), \texttt{entry}(файл, который нужно удалить). Возвращет \texttt{nil}.
	\item \texttt{open-file} -- открыть файл как поток для чтения и записи. Аргументы: \texttt{file}(файл, который нужно открыть). Возвращает объект файла.
	\item \texttt{rename} -- переименовать файл или каталог. При необходимости удалить старые записи длинного имени, а также запись в каталоге, и создать новые записи длинного имени и новую запись в каталоге. Аргументы: \texttt{entry}(файл или каталог, который нужно переименовать), \texttt{new-name}(новое имя файла или каталога). Возвращает \texttt{nil}.
	\item \texttt{free-space} -- получить объём свободного места в текущем разделе диска. Аргументов нет. Возвращает размер в байтах.
	\item \texttt{set-attr} -- изменить атрибуты файла или каталога. Аргументы: \texttt{entry}(файл или каталог, чьи атрибуты необходимо поменять), \texttt{new-attributes}(новые атрибуты файла или каталога). Возвращает \texttt{nil}.
	\item \texttt{fstat} -- получить метаданные файла или каталога. Аргументы: \texttt{entry}(файл или каталог, чьи метаданные надо получить), \texttt{name}(имя файла, чьи метаданные надо получить). Возвращает структуру \texttt{stat}. Структура \texttt{stat} состоит из:
	
		\begin{enumerate}
		\item \texttt{name} -- название файла.
		\item \texttt{size} -- размер файла в байтах. Для каталогов равно 0.
		\item \texttt{create-date} -- дата создания файла.
		\item \texttt{create-time} -- время создания файла.
		\item \texttt{modify-date} -- дата последнего изменения файла.
		\item \texttt{modify-time} -- время последнего изменения файла.
		\item \texttt{access-date} -- дата последнего доступа файла.
		\item \texttt{access-date} -- всегда имеет значение равное \texttt{nil}.
		\item \texttt{isdir} -- флаг, показывающий, является ли файл каталогом.
		\item \texttt{create-time-ms} -- сотые доли секунды создания файла или каталога(1 байт). Имеет значение от 0 до 199(от 0 до 1.99 секунды с шагом 0.01 секунда).
		\item \texttt{attribute} -- атрибуты файла.
		\end{enumerate}
		
	\item \texttt{new-block} -- добавить новый блок к файлу. Аргументы: \texttt{file-object}(объект файла, к которому необходимо добавить новый блок). Возвращает \texttt{nil}.
		
	\end{enumerate}

\item \texttt{FAT32File} -- класс файла файловой системы FAT-32. Наследник класса \texttt{File}. Поля класса \texttt{FAT32File}:

	\begin{enumerate}
	\item \texttt{name} -- имя файла.
	\item \texttt{size} -- размер файла в байтах.
	\item \texttt{position} -- позиция в потоке чтения/записи файла(пара из номера текущего блока и позиции в нём).
	\item \texttt{blocks} -- список номеров блока файла.
	\item \texttt{attributes} -- флаги файла.
	\end{enumerate}
	
Методы класса \texttt{FAT32File}:

	\begin{enumerate}
	\item \texttt{close-file} -- закрыть файл. Аргументы: \texttt{file-object}(объект файла, который необходимо закрыть). Возвращает \texttt{nil}.
	\item \texttt{read-file} -- прочитать определённое количество байт из потока файла, и переместить позицию внутри потока на то же количество. При необходимости обновить номер и содержимое текущего блока. Изменить дату последнего доступа. Аргументы: \texttt{file-object}(объект файла, из которого читать байты), \texttt{size}(количество байт, которое необходимо прочитать). Возвращает массив прочитанных байт.
	\item \texttt{write-file} -- записать массив байт в поток файла, и переместить позицию внутри потока на то же количество. При необходимости выделить новые блоки. При необходимости обновить номер и содержимое текущего блока. Изменить дату и время последнего изменения. Аргументы: \texttt{file-object}(объект файла, в который необходимо записать массив байт), \texttt{buffer}(массив байт, который необходимо записать). Возвращает количество байт, которые было записано.
	\item \texttt{tell-file} -- получить позицию внутри потока файла. Аргументы: \texttt{file-object}(объект файла, позицию в котором необходимо вернуть).
	\item \texttt{seek-file} -- переместить позицию внутри потока файла. При необходимости обновить номер и содержимое текущего блока. Аргументы: \texttt{file-object}(объект файла, позицию в котором необходимо переместить), \texttt{offset}(насколько сместить позицию), \texttt{origin}(откуда смещать позицию). Возвращает значение новой позиции в потоке.	
	\end{enumerate}

\end{enumerate}

\subsection{Модуль файловой системы ISO-9660(CDFS)}

Описание файловой системы ISO-9660(CDFS) приведено выше.

При загрузке файловой системы необходимо прочитать Основной дескриптор тома, запись в каталоге для корневого каталога из основного дескриптора тома.

Из основного дескриптора тома необходимо достать следующую информацию:

\begin{itemize}
\item количество блоков раздела;
\item размер одного блока;
\item размер таблицы путей в байтах;
\item номер сектора, в котором находится таблица путей в формате Little-Endian;
\item запись в каталоге для корневого каталога.
\end{itemize}

Из записи в каталоге для корневого каталога необходимо достать следующую информацию:
\begin{itemize}
\item номер блока с содержимым(только часть с Little-Endian);
\item размер содержимого в байтах(только часть с Little-Endian);
\item дату и время записи;
\item атрибуты корневого каталога.
\end{itemize}

При вызове функции открытия файла, создаётся экземпляр класса CDFSFile. Список блоков файла сохраняется в поле \texttt{blocks}. Во время чтения, позиция внутри текущего блока сохраняется в поле \texttt{position}, номер текущего блока сохраняется в поле \texttt{position}.

Модуль файловой системы ISO-9660(CDFS) включает следующие классы:

\begin{enumerate}
\item \texttt{CDFSFileSystem} -- класс файловой системы ISO-9660(CDFS). Наследник класса \texttt{FileSystem}. Поля класса \texttt{CDFSFileSystem}:

	\begin{enumerate}
	\item \texttt{root-entry} -- запись в каталоге корневого каталога(номер блока с содержимым, размер содержимого, дата и время записи, атрибуты).
	\end{enumerate}
	
Методы класса \texttt{CDFSFileSystem}:

\begin{enumerate}
	\item \texttt{fs-init} -- загрузка файловой системы. Чтение основного дескриптора тома, записи в каталоге для корневого каталога из основного дескриптора тома. Аргументы: \texttt{disk}(номер диска, с которого загружают), \texttt{start-sector}(номер сектора начала раздела), \texttt{sector-count}(количество секторов раздела). Возвращает \texttt{nil}.
	\item \texttt{create-dir} -- файловая система ISO-9660(CDFS) только для чтения, поэтому запись невозможна. Аргументов нет. Вызывает ошибку, что файловая система ISO-9660(CDFS) только для чтения.
	\item \texttt{remove-dir} -- файловая система ISO-9660(CDFS) только для чтения, поэтому запись невозможна. Аргументов нет. Вызывает ошибку, что файловая система ISO-9660(CDFS) только для чтения.
	\item \texttt{load-dir} -- раскрыть и сохранить содержимое каталога. Аргументы: \texttt{dir}(каталог, чьё содержимое нужно загрузить). Возвращает \texttt{nil}.
	\item \texttt{create-file} -- файловая система ISO-9660(CDFS) только для чтения, поэтому запись невозможна. Аргументов нет. Вызывает ошибку, что файловая система ISO-9660(CDFS) только для чтения.
	\item \texttt{remove-file} -- файловая система ISO-9660(CDFS) только для чтения, поэтому запись невозможна. Аргументов нет. Вызывает ошибку, что файловая система ISO-9660(CDFS) только для чтения.
	\item \texttt{open-file} -- открыть файл как поток для чтения и записи. Аргументы: \texttt{file}(файл, который нужно открыть). Возвращает объект файла.
	\item \texttt{rename} -- файловая система ISO-9660(CDFS) только для чтения, поэтому запись невозможна. Аргументов нет. Вызывает ошибку, что файловая система ISO-9660(CDFS) только для чтения.
	\item \texttt{free-space} -- файловая система ISO-9660(CDFS) только для чтения, поэтому запись невозможна. Аргументов нет. Вызывает ошибку, что файловая система ISO-9660(CDFS) только для чтения.
	\item \texttt{set-attr} -- файловая система ISO-9660(CDFS) только для чтения, поэтому запись невозможна. Аргументов нет. Вызывает ошибку, что файловая система ISO-9660(CDFS) только для чтения.
	\item \texttt{fstat} -- получить метаданные файла или каталога. Аргументы: \texttt{dir}(каталог с файлом, чьи метаданные надо получить), \texttt{name}(имя файла, чьи метаданные надо получить). Возвращает структуру \texttt{stat}. Структура \texttt{stat} состоит из:
	
	\begin{enumerate}
		\item \texttt{name} -- название файла.
		\item \texttt{size} -- размер файла в байтах. Для каталогов равно 0.
		\item \texttt{create-date} -- дата создания файла.
		\item \texttt{create-time} -- время создания файла.
		\item \texttt{modify-date} -- всегда имеет значение равное \texttt{nil}.
		\item \texttt{modify-time} -- всегда имеет значение равное \texttt{nil}.
		\item \texttt{access-date} -- всегда имеет значение равное \texttt{nil}.
		\item \texttt{access-date} -- всегда имеет значение равное \texttt{nil}.
		\item \texttt{isdir} -- флаг, показывающий, является ли файл каталогом.
		\item \texttt{flags} -- атрибуты файла.
	\end{enumerate}
	
	\item \texttt{new-block} -- добавить новый блок к файлу. Аргументы: \texttt{file-object}(объект файла, к которому необходимо добавить новый блок). Вызывает ошибку, что файловая система ISO-9660(CDFS) только для чтения.
	
\end{enumerate}

\item \texttt{CDFSFile} -- класс файла файловой системы ISO-9660(CDFS). Наследник класса \texttt{File}. Поля класса \texttt{CDFSFile}:

	\begin{enumerate}
	\item \texttt{name} -- имя файла.
	\item \texttt{size} -- размер файла в байтах.
	\item \texttt{position} -- позиция в потоке чтения/записи файла(пара из номера текущего блока и позиции в нём).
	\item \texttt{blocks} -- список номеров блока файла.
	\end{enumerate}

Методы класса \texttt{CDFSFile}:

	\begin{enumerate}
	\item \texttt{close-file} -- закрыть файл. Аргументы: \texttt{file-object}(объект файла, который необходимо закрыть). Возвращает \texttt{nil}.
	\item \texttt{read-file} -- прочитать определённое количество байт из потока файла, и переместить позицию внутри потока на то же количество. При необходимости обновить номер и содержимое текущего блока. Изменить дату последнего доступа. Аргументы: \texttt{file-object}(объект файла, из которого читать байты), \texttt{size}(количество байт, которое необходимо прочитать). Возвращает массив прочитанных байт.
	\item \texttt{write-file} -- файловая система ISO-9660(CDFS) только для чтения, поэтому запись невозможна. Аргументов нет. Вызывает ошибку, что файловая система ISO-9660(CDFS) только для чтения.
	\item \texttt{tell-file} -- получить позицию внутри потока файла. Аргументы: \texttt{file-object}(объект файла, позицию в котором необходимо вернуть).
	\item \texttt{seek-file} -- переместить позицию внутри потока файла. При необходимости обновить номер и содержимое текущего блока. Аргументы: \texttt{file-object}(объект файла, позицию в котором необходимо переместить), \texttt{offset}(насколько сместить позицию), \texttt{origin}(откуда смещать позицию). Возвращает значение новой позиции в потоке.
	\end{enumerate}

\end{enumerate}

\subsection{Модуль работы с блоками диска}

Блок -- один или несколько секторов подряд.

Из-за того, что размер блока у разных файловых систем может быть разным(он определяется при загрузке файловой системы), необходима глобальная переменная \texttt{block-sectors}, которая будет хранить количество секторов на один блок.

Также из-за того, что у файловых систем разные номера первого блока области данных, необходима глобальная переменная \texttt{block-offset}, которая будет хранить смещение в секторах от начала раздела для корректной нумерации блоков.

Из-за того, что номер первого блока области данных у разных файловых систем имеет разное значение, необходима глобальная переменная \texttt{blocks-start-num}, которая будет хранить минимальный номер блока области данных, чтобы при возникновении ошибок(в качестве номера блока для чтения или записи попадает номер, который указывает на блок не в области данных) не перезаписать часть структуры файловых систем.

Модуль работы с блоками диска включает следующие переменные:

\begin{enumerate}
	\item \texttt{disk-num} -- номер диска, с которого читают блоки.
	\item \texttt{block-offset} -- смещение в секторах для отсчёта блоков.
	\item \texttt{block-sectors} -- количество секторов на один блок.
	\item \texttt{blocks-start-num} -- минимальный номер блока для чтения или записи.
\end{enumerate}

Модуль работы с блоками диска включает следующие функции:

\begin{enumerate}
	\item \texttt{set-block-disk} -- установить диск, с которого читать блоки. Аргументы: \texttt{disk}(номер диска, с которого читать блоки). Возвращает \texttt{nil}.
	\item \texttt{set-block-size} -- установить размер блока диска. Аргументы: \texttt{size}(размер блока в секторах). Возвращает \texttt{nil}.
	\item \texttt{set-block-offset} -- установить смещение в секторах для отсчёта блоков диска. Аргументы: \texttt{offset}(смещение в секторах). Возвращает \texttt{nil}.
	\item \texttt{set-blocks-start-num} -- установить минимальный номер блока для чтения или записи. Аргументы: \texttt{blocks-start-num}(минимальный номер блока для чтения или записи). Возвращает \texttt{nil}.
	\item \texttt{block-read} -- прочитать блок диска. Аргументы: \texttt{num}(номер блока, который читают). Возвращает массив прочитанных байт.
	\item \texttt{block-write} -- записать в блок диска. Аргументы: \texttt{num}(номер блока, в который записывают), \texttt{buf}(массив байт). Возвращает \texttt{nil}.
	\item \texttt{get-blocks-pos} -- получить номер блока и смещение по позиции в списке блоков. Аргументы: \texttt{blocks}(список блоков, в котором ищут позицию), \texttt{pos}(смещение в байтах). Возвращает пару из номера блока и позиции в нём.
	\item \texttt{get-offset-from-pos} -- получить смещение в байтах из позиции. Аргументы: \texttt{pos}(индекс блока и смещение в блоке). Возвращает смещение в байтах.
\end{enumerate}

\subsection{Модуль драйвера ATA}

Интерфейс ATA используется для работы с секторами диска.

Сектор диска -- минимальная единица хранения информации на диске.

Модуль драйвера ATA включает следующие функции:

\begin{enumerate}
	\item \texttt{ata-get-status-reg} -- получение значения бита регистра статуса. Аргументы: \texttt{bit}(номер бита регистра статуса). Возвращает значение бита регистра статуса.
	\item \texttt{ata-set-dev} -- установить номер устройства. Аргументы: \texttt{dev}(номер устройства). Возвращает \texttt{nil}.
	\item \texttt{ata-set-lba} -- установить номер сектора LBA. Аргументы: \texttt{sec}(номер первого сектора), \texttt{num}(количество секторов). Возвращает \texttt{nil}.
	\item \texttt{ata-check-error} -- проверить на ошибку. Аргументов нет. Возвращает \texttt{t}, если есть ошибка, иначе \texttt{nil}.
	\item \texttt{ata-set-command} -- установить команду контроллера. Аргументы: cmd(команда контроллера). Возвращает \texttt{nil}.
	\item \texttt{ata-read} -- прочитать данные контроллера. Аргументы: \texttt{size}(количество байт). Возвращает массив прочитанных байт.
	\item \texttt{ata-write} -- записать массив в контроллер. Аргументы: \texttt{arr}(массив байт). Возвращает \texttt{nil}.
	\item \texttt{ata-identify} -- прочитать служебную информацию. Аргументов нет. Возвращает массив со служебной информацией: число головок, цилиндров, секторов, серийный номер и т.д.
	\item \texttt{ata-read-sectors} -- прочитать с секторов диска. Аргументы: \texttt{dev}(номер устройства), \texttt{start}(номер первого сектора), \texttt{num}(количество секторов). Возвращает массив прочитанных байт.
	\item \texttt{ata-write-sectors} -- записать массив байт на сектора диска. Аргументы: \texttt{dev}(номер устройства), \texttt{start}(номер первого сектора), \texttt{num}(количество секторов), \texttt{arr}(массив байт). Возвращает \texttt{nil}.
\end{enumerate}